% Generated draft: requires mathematical and source review.
\documentclass[11pt]{article}
\usepackage[margin=25mm]{geometry}
\usepackage{fontspec}
\setmainfont{Latin Modern Roman}
\usepackage{amsmath,amssymb,mathtools,unicode-math}
\setmathfont{Latin Modern Math}
\usepackage{enumitem,needspace,xurl,hyperref}
\setlength{\parindent}{0pt}
\setlength{\parskip}{5pt}
\setlength{\emergencystretch}{3em}
\newcommand{\N}{\mathbb{N}}
\newcommand{\Z}{\mathbb{Z}}
\newcommand{\Q}{\mathbb{Q}}
\newcommand{\R}{\mathbb{R}}
\newcommand{\C}{\mathbb{C}}
\newcommand{\F}{\mathbb{F}}
\newcommand{\A}{\mathbb{A}}
\newcommand{\PP}{\mathbb{P}}
\newcommand{\E}{\mathbb{E}}
\newcommand{\eps}{\varepsilon}
\newcommand{\dd}{\,\mathrm{d}}
\begin{document}
\section*{UnsolvedMath ID: 30004167}
\textbf{Source problem code:} catalogue entry
\subsection{Definitions and mathematical statement}
The supplied record concerns topological periodic cyclic homology, denoted by \(TP\), with \(TP_k\) denoting its degree-\(k\) group. It uses relative expressions \(TP_k(R/S_n)\) and \(TP_k(R/S_n^{\otimes 2})\), but does not define \(S_n\), the parameter \(n\), or the precise relative construction.

Let \(R\) be an \(S_n\)-algebra in the intended setup. The proposed identification concerns the groups
\[
\operatorname{Ext}_{TP_0(R/S_n^{\otimes 2})}^{j}
\bigl(TP_0(R/S_n),TP_i(R/S_n)\bigr).
\]
Here \(\operatorname{Ext}_A^j(M,N)\) denotes an Ext group over the indicated coefficient ring \(A\), once the intended module category and module structures on \(M\) and \(N\) have been specified. The question is whether these groups identify with the appropriate prismatic cohomology groups of \(R\), and whether the associated Adams-type spectral sequence displayed in the source,
\[
\operatorname{Ext}_{TP_0(R/S_n^{\otimes 2})}^{j}
\bigl(TP_0(R/S_n),TP_i(R/S_n)\bigr)
\Longrightarrow TP_{i-j}(R),
\]
is the Bhatt--Morrow--Scholze spectral sequence.

\emph{Unresolved qualifications.} The supplied record does not specify the allowed \(S_n\)-algebras, the ranges of \(n,i,j\), the tensor-product base or whether the tensor product is derived, or the coefficient actions needed to form the Ext groups. It also does not specify a prime, completion conventions, a prism or other base for prismatic cohomology, the cohomological grading and possible twists in the proposed identification, or the page, filtration, and convergence conditions of the displayed spectral sequence. Thus the display preserves the supplied target, but does not by itself determine a fully specified conjecture or assert the existence or convergence of a spectral sequence. These missing choices require clarification from the underlying setup.
\subsection{Short English statement}
For an \(S_n\)-algebra \(R\), the supplied conjectural question asks whether the displayed Ext groups compute prismatic cohomology and whether the associated Adams-type spectral sequence is the Bhatt--Morrow--Scholze spectral sequence. The base, admissible algebras, coefficient structures, gradings, and convergence conventions remain unspecified in the supplied record.
\subsection{Sources}
\begin{itemize}
\item Local registry reference: supplied UnsolvedMath source record, permanent numeric ID \(30004167\), exact source problem code \texttt{catalogue entry}, in the Oberwolfach Reports Open Problems collection. The registry title is ``Prismatic Cohomology from Topological Cyclic Homology''; the displayed statement specifically uses topological periodic cyclic homology.
\item Supplied report citation: \emph{Algebraic K-theory (2020)}, Oberwolfach Reports, pp.~1737--1790, DOI \texttt{10.4171/owr/2019/29}. Supplied source URL and literature-source link, labeled ``OWR primary report'': \url{https://doi.org/10.4171/owr/2019/29}. The record describes Oberwolfach Reports as documenting mathematical workshops and their open-problem discussions. It records the proposed year as 2019 and supplies no individual proposer.
\item Supplied verification metadata labels the statement \texttt{corrected\_verified} and reports that it was checked against the official Oberwolfach report, preserving its hypotheses, notation, and open mathematical target. This is a claim in the supplied curation, not an independent verification for this draft.
\item Supplied literature assessment, checked on \texttt{2026-08-22}: the record labels the problem open, describes the proposed Ext and spectral-sequence identifications as conjectural, and reports that no primary source located settles this exact construction in the stated generality. This is supplied dated curation derived from the OWR record, not independently verified. The record requests independent verification of time-sensitive status and cited claims.
\end{itemize}
\end{document}
